\chapter{Conclusão}
\label{ch:consideracoes}

%%%%%%%%%%%%%%%%%%%%%%%%%%%%%%%%%%%%%%%%%%%%%%%%%%%%%%%%%%%

\section{Conclusões}
\label{sec:conclusoes}

Escrever\ldots

%%%%%%%%%%%%%%%%%%%%%%%%%%%%%%%%%%%%%%%%%%%%%%%%%%%%%%%%%%%

\section{Contribuições}
\label{sec:conclusao}

Este trabalho propôs um modelo para suportar o gerenciamento de perfis em ambientes ubíquos, permitindo a inferência de trilhas de entidades através da definição de regras pelas aplicações. Embora o eProfile tinha sido integrado com o LOCAL e o UniManager no cenário apresentado no Capítulo~\ref{ch:avaliacao}, sua proposta é genérica e pode ser utilizada para qualquer aplicação ou conjunto de aplicações que utilize perfis para aperfeiçoamento da personalização. O uso de perfis automáticos permite que sistemas ubíquos sejam mais eficientes na compreensão da intenção do usuário, pois isso significa que sistemas deste tipo podem utilizar informações mais precisas, atualizadas e amplas para suprir as necessidades das aplicações.

O objetivo do cenário foi demonstrar que é viável integrar sistemas ubíquos com o eProfile, enriquecendo a geração de perfis destas aplicações e permitindo que elas se adaptem aos seus usuários. O cenário demonstra as quatro contribuições do modelo.

A primeira contribuição é o uso de trilhas para extração de perfis. A utilização de trilhas enriquece o banco de dados de informações que pode ser utilizado para a geração de perfis ao oferecer uma visão histórica das ações das entidades dentro dos contextos. No cenário apresentado, decisões passadas a respeito de assuntos de trabalhos e debates foram utilizadas para extrair as preferências que compuseram os perfis dos estudantes.

A segunda contribuição é a geração de perfis dinâmicos. Ao atualizar os perfis com novas informações presentes nas trilhas, é possível manter o perfil sincronizado com o que está ocorrendo com a entidade. Isto permite que as aplicações personalizem a experiência do usuário. No cenário apresentado, os perfis são atualizados toda vez que o estudante toma uma decisão referente a um debate ou trabalho.

A terceira contribuição é a capacidade de gerenciar regras de inferência dinâmicas e modelos dinâmicos. Ao permitir que as regras e os modelos sejam atualizados a qualquer momento, o eProfile facilita a adaptação da aplicação ao seu usuário. Um segundo experimento foi realizado, onde as regras de inferência para a geração dos perfis foram alteradas durante a execução da simulação, resultando numa mudança positiva no resultado do cenário.

A quarta contribuição é a operabilidade semântica do modelo. Isto abre a possibilidade de uma aplicação utilizar informações geradas por uma aplicação diferente para gerar os perfis de suas entidades. Isto é demonstrado no cenário onde o UniManager utiliza as informações geradas pelo LOCAL (as disciplinas opcionais escolhidas pelos estudantes) como parte dos dados usados para criar os perfis utilizados pelo UniManager. Esta informação depois é utilizada para sugerir disciplinas opcionais aos estudantes.

\section{Trabalhos Futuros}

O eProfile é uma proposta inicial, tal como a integração com o LOCAL e o UniManager. De acordo com os resultados obtidos na validação do modelo, apresentam-se algumas proposta para dar continuidade ao desenvolvimento e ampliação do modelo.

Existe espaço para ampliar as avaliações do eProfile. As funcionalidades do eProfile foram validadas através da simulação de um cenário. Uma avaliação ideal envolveria a geração de perfis por entidades reais, ao invés da simulação, onde o modelo poderia ser posto à prova em um ambiente real, durante um tempo extenso. Embora a simulação traga as vantagens de várias repetições em um curto espaço de tempo, uma avaliação em um ambiente real traria um resultado mais preciso, e levaria em conta aspectos não previstos anteriormente.

A simulação utilizou um conjunto pequeno de entidades, e não houve avaliação de desempenho. Uma vez que o modelo é genérico, outros ambientes poderiam utilizar um numero muito maior de entidades e necessitar de perfis gerados em um período de tempo mais curto. Desta forma, a avaliação do desempenho da execução do eProfile sobre um grande número de trilhas em um tempo curto poderia apontar oportunidades de melhora de desempenho no modelo.

O cenário citou a integração do eProfile a duas aplicações (LOCAL e UniManager). Cada aplicação apresenta características próprias, e a integração do eProfile a outras aplicações poderá apresentar oportunidades para expansão das funcionalidades. Além disso, a integração com ambas aplicações ocorreu dentro de um contexto educacional. A integração com aplicações de outros contextos pode também indicar possibilidades para implementação de novas funcionalidades.

Como foi mencionado, o modelo é uma proposta inicial, e existem várias oportunidades para ampliação. Entre elas estão a compatibilização com outros gerenciadores de trilhas, expansão das ontologias básicas, suporte à operação desconectada e/ou instável, descentralização do servidor, suporte a protocolos de segurança e confidencialidade, entre outros.

%(i) a integração de mais aplicações com o eProfile, de modo a testar e expandir as capacidades do eProfile, (ii) o desempenho de uma grande quantidade de trilhas precisa ser avaliada, uma vez que o experimento considerou um pequeno número de interações e (iii) a simulação ideal envolveria trilhas geradas por entidades reais, ao invés de uma simulação, 
